%=====================================================================
\chapter{Decomposition, and Denormalising on Purpose}\label{ch:decomposition}
%=====================================================================
\locator{5}{Part V --- Getting the Design Right}

\begin{outcomes}
By the end of this chapter you will be able to:
\begin{tight}
  \item \textbf{Define} a lossless-join decomposition and \textbf{apply} the
        one-line test for it.
  \item \textbf{Show} what a lossy decomposition does to a query, by counting
        its spurious tuples.
  \item \textbf{Define} dependency preservation and \textbf{identify} a
        decomposition that loses a dependency.
  \item \textbf{Decide} between BCNF and 3NF when they conflict.
  \item \textbf{Measure} a denormalisation --- what it saves per read, what it
        costs per write --- and \textbf{compute} where it breaks even.
\end{tight}
\end{outcomes}

\begin{prereq}
\chref{normalforms}: the normal forms, and its measurement of the wide table.
\chref{fds} for closure, which the lossless test uses. You need
\texttt{code/m19\_decomposition.py}.
\end{prereq}

\begin{keyterms}
decomposition $\bullet$ lossless join $\bullet$ lossy join $\bullet$ spurious
tuple $\bullet$ dependency preservation $\bullet$ projection of dependencies
$\bullet$ 3NF synthesis $\bullet$ denormalisation $\bullet$ summary table
$\bullet$ materialised data $\bullet$ trigger $\bullet$ break-even $\bullet$
read--write ratio $\bullet$ staleness
\end{keyterms}

\section{Two Things a Split Must Not Lose}

\chref{normalforms} decomposed \textit{report} into six tables by splitting on
violating dependencies. Any split \emph{can} be made; not every split is
correct. A decomposition must keep two things, and they are independent.

\begin{definitionbox}[Lossless join]
A decomposition of $R$ into $R_1$ and $R_2$ is \term{lossless} if joining them
back gives exactly $R$ --- no tuple missing, \textbf{no tuple added}.

\smallskip
\textbf{The test.} It is lossless if and only if the attributes they share
determine all of one side:
\[
  (R_1 \cap R_2) \rightarrow R_1 \quad\text{or}\quad
  (R_1 \cap R_2) \rightarrow R_2 .
\]
In words: the shared attributes are a superkey of one of the halves. Compute
$(R_1 \cap R_2)^{+}$, as in \chref{fds}, and look.

\smallskip
A \term{lossy} decomposition never loses tuples --- it \emph{adds}
\term{spurious tuples}, rows that look like facts and were never true.
\end{definitionbox}

\begin{definitionbox}[Dependency preservation]
A decomposition \term{preserves dependencies} if every dependency of $R$ can be
checked within a single one of the $R_i$ --- or follows from dependencies that
can.

\smallskip
It matters because a database enforces rules one table at a time. A
dependency that spans two tables can only be checked by joining them, and
nothing in the schema does that on every insert.
\end{definitionbox}

\begin{measured}[A split that invents facts, a split that loses a rule, and a summary table priced]
\textsf{\footnotesize \texttt{code/m19\_decomposition.py}, seeded database
(4{,}000 students), SQLite 3.51.2, Intel i5-6600, Windows 11.}

\medskip
\textbf{1. Lossy against lossless}

\begin{tabular}{@{}L{9.6cm}R{2.8cm}@{}}
\multicolumn{2}{@{}l}{\textbf{lossy:} $R_1$(roll\_no, course\_code) and
  $R_2$(course\_code, grade)} \\
graded (roll\_no, course\_code, grade) facts & 39{,}152 \\
$R_1 \bowtie R_2$ on course\_code & 182{,}920 \\
\textbf{spurious tuples --- never true} & \textbf{143{,}768} \\
\multicolumn{2}{@{}l}{\textbf{lossless:} $S_1$(roll\_no, name, programme) and
  $S_2$(roll\_no, semester, cnic)} \\
student rows, and $S_1 \bowtie S_2$ & 4{,}000 and 4{,}000 \\
\end{tabular}

\smallskip
The lossy join returned \textbf{4.7$\times$} the real facts: every student now
appears to hold every grade that anybody got in the course. The test explains
it in one line --- $R_1 \cap R_2 = \{$course\_code$\}$, a key of neither half;
$S_1 \cap S_2 = \{$roll\_no$\}$, the key of both.

\medskip
\textbf{2. Lossless, but a dependency is lost}

Two tables, \textit{course\_dept}(course\_code, dept\_code) and
\textit{course\_dname}(course\_code, dept\_name), both keyed on course\_code ---
so the join is lossless. Two courses, both CS. A second name for CS:

\begin{tabular}{@{}L{9.6cm}L{2.8cm}@{}}
insert \texttt{('CS-150', 'Computing')} & \textbf{accepted} \\
the join: CS $\rightarrow$ & \texttt{'Computer Science'} \\
the join: CS $\rightarrow$ & \texttt{'Computing'} \\
\end{tabular}

\smallskip
\fd{dept\_code}{dept\_name} is violated, and nothing could refuse it: the rule
spans two tables, and a \sqlkw{CHECK} sees one row of one table. This is
exactly the decomposition \chref{normalforms} warned the wrong splitting order
produces.

\medskip
\textbf{3. Denormalisation, priced} --- enrolments and credit hours per
department

\begin{tabular}{@{}L{8.4cm}R{2.6cm}@{}}
read: computed by two joins & 69.79\,ms \\
read: from a summary table & 0.060\,ms \\
\quad \textbf{saved per read} & \textbf{69.73\,ms} \\
write: one enrolment, no trigger & 28.464\,ms \\
write: one enrolment, trigger maintaining the summary & 32.143\,ms \\
\quad \textbf{added per write} & \textbf{3.679\,ms} \\
summary agreed with the joins, before and after 400 writes & yes, yes \\
\end{tabular}

\smallskip
\textbf{Break-even: one read every 19 writes.} Above that ratio the summary
table is cheaper overall; below it, computing on demand is. A department report
is read far more often than once per nineteen enrolments, so for this report
the summary table wins --- by a margin that is not close.

\smallskip
Most of the 28\,ms per write is the \sqlkw{COMMIT} --- making one enrolment
durable --- and the trigger's two \sqlkw{UPDATE}s ride inside that same
transaction. That is why the trigger adds 13\%, not 200\%.
\end{measured}

\dsfig{%
\begin{tikzpicture}[
  cell/.style={draw=primary!35, line width=0.5pt, minimum width=1.25cm,
               minimum height=0.40cm, font=\scriptsize\ttfamily,
               text=primary!65!black, inner sep=0pt},
  hd/.style={cell, fill=primary, text=white, font=\scriptsize\bfseries},
  bad/.style={cell, fill=accent!14, text=accent},
  ttl/.style={font=\scriptsize\bfseries, text=#1, anchor=south},
  nt/.style={font=\scriptsize\itshape, text=neutral, align=center}
]
  % original
  \node[ttl=primary] at (-5.8,0.5) {the facts};
  \foreach \c [count=\i] in {roll, course, grade}{
    \node[hd] at (-7.05+\i*1.25-1.25+1.25,0.2) {\c};}
  \foreach \r/\a/\b/\g in {1/Ayesha/CS-103/A, 2/Bilal/CS-103/C}{
    \node[cell] at (-5.80,-0.4*\r+0.2) {\a};
    \node[cell] at (-4.55,-0.4*\r+0.2) {\b};
    \node[cell] at (-3.30,-0.4*\r+0.2) {\g};}
  \node[nt] at (-4.55,-1.05) {2 rows};

  % R1, R2
  \node[ttl=secondary!70!black] at (-0.6,1.65) {$R_1$};
  \node[hd] at (-1.25,1.35) {roll};  \node[hd] at (0.0,1.35) {course};
  \node[cell] at (-1.25,0.95) {Ayesha}; \node[cell] at (0.0,0.95) {CS-103};
  \node[cell] at (-1.25,0.55) {Bilal};  \node[cell] at (0.0,0.55) {CS-103};
  \node[ttl=secondary!70!black] at (-0.6,-0.35) {$R_2$};
  \node[hd] at (-1.25,-0.65) {course}; \node[hd] at (0.0,-0.65) {grade};
  \node[cell] at (-1.25,-1.05) {CS-103}; \node[cell] at (0.0,-1.05) {A};
  \node[cell] at (-1.25,-1.45) {CS-103}; \node[cell] at (0.0,-1.45) {C};

  \draw[-{Stealth[length=2mm]}, primary!50, line width=0.9pt]
       (-2.55,-0.2) -- (-2.0,-0.2);
  \draw[-{Stealth[length=2mm]}, primary!50, line width=0.9pt]
       (0.75,-0.2) -- (1.35,-0.2) node[midway, above,
       font=\scriptsize\bfseries, text=primary] {$\bowtie$};

  % rejoined
  \node[ttl=accent] at (3.35,0.5) {$R_1 \bowtie R_2$ on course};
  \foreach \c [count=\i] in {roll, course, grade}{
    \node[hd] at (1.6+\i*1.25-1.25+0.5,0.2) {\c};}
  \foreach \r/\a/\g/\s in {1/Ayesha/A/cell, 2/Ayesha/C/bad,
                           3/Bilal/A/bad, 4/Bilal/C/cell}{
    \node[\s] at (2.10,-0.4*\r+0.2) {\a};
    \node[\s] at (3.35,-0.4*\r+0.2) {CS-103};
    \node[\s] at (4.60,-0.4*\r+0.2) {\g};}
  \node[nt, text width=3.9cm, text=accent] at (3.35,-1.95)
       {4 rows: \textbf{2 never true}. Ayesha did not get a C; Bilal did not
        get an A};
\end{tikzpicture}}%
{A lossy split on the smallest possible scale. The two halves share only
\textit{course}, which is a key of neither, so the join pairs every student in
the course with every grade in it. On the seeded data this is 143{,}768
invented facts.}{lossy}

\begin{alertbox}[What this does and does not license]
\begin{tight}
  \item \textbf{Always test a decomposition for losslessness.} A lossy one
        produces 4.7 false facts for every true one here, and no query
        afterwards can tell them apart.
  \item \textbf{The summary table is a denormalisation, and it is correct only
        while the trigger runs.} Drop the trigger, insert through a path that
        bypasses it, or restore an old backup of one table, and the summary is
        silently wrong. That risk is what the 69\,ms buys.
  \item \textbf{Denormalise by adding, never by replacing.} The normalised
        tables stayed. The summary is derived from them and can be rebuilt
        from them at any moment --- which is the difference between this and
        \chref{normalforms}'s \textit{report}, where the wide table was the
        only copy.
  \item \textbf{Decide with a ratio, not a feeling.} One read per nineteen
        writes is a number; ``reports are slow'' is not.
\end{tight}
\end{alertbox}

\section{When BCNF and Dependency Preservation Conflict}

\begin{conceptbox}[The trade, and the rule for making it]
Every relation has a \textbf{lossless, dependency-preserving decomposition into
3NF}. The \term{3NF synthesis} algorithm finds one: take a minimal cover
(\chref{fds}), make one relation per dependency, add one relation holding a
candidate key if none already does.

\smallskip
A lossless BCNF decomposition always exists too --- but it is \emph{not}
always dependency-preserving. \chref{normalforms}'s
\relation{teaches}{student, course, instructor} is the case: BCNF forces
\fd{student, course}{instructor} out of every single table.

\smallskip
\textbf{The rule:} prefer BCNF; when it would lose a dependency you need
enforced, stop at 3NF and accept the small redundancy that remains. A
redundancy you can see is safer than a rule nothing enforces. For the
university schema the question never arises --- the six tables are in BCNF
\emph{and} preserve every dependency.
\end{conceptbox}

\begin{worked}[Testing two decompositions of \textit{section}]
\relation{section}{\relkey{course\_code}, \relkey{semester},
\relkey{section\_no}, room, emp\_id}, with
\fd{course\_code, semester, section\_no}{room, emp\_id}.

\medskip
\textbf{Decomposition A:} $A_1$(course\_code, semester, section\_no, room)
and $A_2$(course\_code, semester, section\_no, emp\_id).

\smallskip
$A_1 \cap A_2 = \{$course\_code, semester, section\_no$\}$, whose closure is
everything: a superkey of both. \textbf{Lossless.} The dependency splits into
\fd{\dots}{room} in $A_1$ and \fd{\dots}{emp\_id} in $A_2$ --- each checkable in
one table. \textbf{Preserving.} Correct, and pointless: it is two tables where
one would do, and every split is a join somebody pays for on every read.

\medskip
\textbf{Decomposition B:} $B_1$(course\_code, room) and $B_2$(room, semester,
section\_no, emp\_id).

\smallskip
$B_1 \cap B_2 = \{$room$\}$. Does \textit{room} determine all of either half?
No: one room hosts many courses and many sections. \textbf{Lossy.} Joined back
on room, every course would appear in every section that ever met in the same
room --- the pattern of part~1 of the measurement, in a different table.

\medskip
\textbf{The lesson.} The test took one closure each time. B ``looks
reasonable'' --- room is a natural thing to split on --- and is wrong in a way
the data would never report.
\end{worked}

\section{Triggers, and the Summary Table Written Out}

\begin{examplebox}[The summary table that was measured]
\begin{lstlisting}[language=sql]
CREATE TABLE dept_stats (
  dept_code   TEXT PRIMARY KEY NOT NULL REFERENCES department(dept_code),
  n_enrol     INTEGER NOT NULL,
  sum_credits INTEGER NOT NULL
) STRICT;

CREATE TRIGGER stats_ins AFTER INSERT ON enrolment BEGIN
  UPDATE dept_stats
     SET n_enrol     = n_enrol + 1,
         sum_credits = sum_credits +
           (SELECT credit_hours FROM course
             WHERE course_code = NEW.course_code)
   WHERE dept_code =
           (SELECT dept_code FROM course
             WHERE course_code = NEW.course_code);
END;
\end{lstlisting}

\smallskip
\texttt{NEW} is the row just inserted. A matching \texttt{stats\_del} trigger
on \sqlkw{DELETE} uses \texttt{OLD}. \textbf{An \sqlkw{UPDATE} trigger is
needed too}, and the harness does not have one: an enrolment moved from one
course to another would leave both departments' figures wrong. That gap is
\chref{views}' starting point, and the reason a view is often the better
answer.
\end{examplebox}

\begin{pitfall}
\begin{tight}
  \item \textbf{Splitting on an attribute that is not a key of either half.}
        Lossy: 143{,}768 spurious tuples here.
  \item \textbf{Expecting a lossy join to lose rows.} It adds them. The answer
        is too big, not too small, and looks plausible.
  \item \textbf{A rule that spans two tables.} No \sqlkw{CHECK} can enforce
        it; it needs a trigger, or a design that keeps it in one table.
  \item \textbf{Denormalising by replacing the normalised tables.} Add the
        summary; keep the source of truth.
  \item \textbf{Forgetting the \sqlkw{UPDATE} trigger.} Insert and delete are
        maintained; a changed row silently is not.
  \item \textbf{Deciding to denormalise without the ratio.} Measure reads and
        writes; compute the break-even.
\end{tight}
\end{pitfall}

\begin{lab}[Invent 143{,}768 facts, then price a summary table]
\begin{lstlisting}[language=bash]
# --- 1. The seeded database ----------------------------------------
python code/schema/seed.py university.db
sqlite3 university.db
\end{lstlisting}
\begin{lstlisting}[language=sql]
-- --- 2. A lossy split -------------------------------------------
CREATE TEMP TABLE r1 AS SELECT DISTINCT roll_no, course_code
                          FROM enrolment WHERE grade IS NOT NULL;
CREATE TEMP TABLE r2 AS SELECT DISTINCT course_code, grade
                          FROM enrolment WHERE grade IS NOT NULL;
SELECT COUNT(*) FROM (SELECT DISTINCT roll_no, course_code, grade
                        FROM enrolment WHERE grade IS NOT NULL);
SELECT COUNT(*) FROM r1 JOIN r2 USING (course_code);
-- How many facts were invented? Apply the one-line test to explain.

-- --- 3. Pick one student and see what the lossy join says ------
SELECT grade FROM r1 JOIN r2 USING (course_code)
 WHERE roll_no = (SELECT MIN(roll_no) FROM r1) LIMIT 10;
SELECT grade FROM enrolment
 WHERE roll_no = (SELECT MIN(roll_no) FROM r1) AND grade IS NOT NULL;

-- --- 4. The summary table and its trigger -----------------------
-- Type in dept_stats and stats_ins from this chapter, then fill it:
INSERT INTO dept_stats
  SELECT d.dept_code, COUNT(e.roll_no), COALESCE(SUM(c.credit_hours),0)
    FROM department d
    LEFT JOIN course c    ON c.dept_code   = d.dept_code
    LEFT JOIN enrolment e ON e.course_code = c.course_code
   GROUP BY d.dept_code;

-- --- 5. Time both reads -----------------------------------------
.timer on
SELECT * FROM dept_stats;
SELECT c.dept_code, COUNT(*), SUM(c.credit_hours)
  FROM enrolment e JOIN course c ON c.course_code = e.course_code
 GROUP BY c.dept_code;

-- --- 6. Break it --------------------------------------------------
UPDATE enrolment SET course_code =
       (SELECT course_code FROM course WHERE dept_code = 'MA' LIMIT 1)
 WHERE rowid = (SELECT MIN(rowid) FROM enrolment);
-- Re-run both reads from step 5. Do they still agree? Why not?
\end{lstlisting}
\begin{lstlisting}[language=bash]
# --- 7. The measured version ----------------------------------------
python code/m19_decomposition.py
\end{lstlisting}

\smallskip
\textbf{What to notice.} Step 6 is the price of the 69\,ms. The summary table is
correct exactly as long as every change to \textit{enrolment} passes through a
trigger that knows about it --- and the harness, like most first attempts,
forgot \sqlkw{UPDATE}. (Step 6 may also be refused if the move would duplicate
an enrolment's key; pick another row if so.)
\end{lab}

\begin{checkpoint}
\begin{tightnum}
  \item Apply the lossless test to $R_1$(roll\_no, course\_code) and
        $R_2$(course\_code, grade), and explain from the result why the join
        produced 4.7 times the real facts.
  \item Give a decomposition of the university schema that is lossless but not
        dependency-preserving, name the dependency it loses, and show an insert
        that violates the dependency without any table refusing it.
  \item The summary table saves 69.73\,ms per read and costs 3.679\,ms per
        write. Show the break-even calculation, and say what would change it
        if the database moved to a fast SSD where a commit took 1\,ms.
\end{tightnum}
\end{checkpoint}

\begin{chaptersummary}
\begin{tight}
  \item A decomposition must be \textbf{lossless} (the join gives back exactly
        $R$) and should be \textbf{dependency-preserving} (every rule checkable
        in one table). The two are independent.
  \item \textbf{Lossless test:} the shared attributes must be a superkey of one
        half. One closure decides it.
  \item \textbf{Measured:} a lossy split of \textit{enrolment} produced
        \textbf{143{,}768 spurious tuples} --- 4.7$\times$ the real facts. A
        lossy join adds rows; it does not lose them.
  \item \textbf{Measured:} a lossless split that lost
        \fd{dept\_code}{dept\_name} \textbf{accepted two names for CS}. A rule
        spanning two tables cannot be enforced by a \sqlkw{CHECK}.
  \item 3NF synthesis is always lossless and preserving; BCNF is always
        lossless but not always preserving. Prefer BCNF; stop at 3NF when BCNF
        would cost a rule you need.
  \item \textbf{Measured:} a trigger-maintained summary table saved
        \textbf{69.73\,ms per read} for \textbf{3.679\,ms per write} --- break
        even at \textbf{one read per 19 writes}.
  \item Most of a write's cost is the \sqlkw{COMMIT}; a trigger rides inside
        it, which is why it added 13\%.
  \item Denormalise by \emph{adding} derived data beside the normalised
        tables, never by replacing them --- and maintain it on insert, delete
        \emph{and} update.
\end{tight}
\end{chaptersummary}

\begin{reviewq}
  \item \textbf{(MCQ)} A lossy decomposition, when joined back, produces:
        (a)~fewer tuples; (b)~the same tuples; (c)~extra, spurious tuples;
        (d)~an error.
  \item \textbf{(MCQ)} $R_1 \cap R_2 = \{A\}$ is lossless if: (a)~$A$ is a
        superkey of $R_1$ or of $R_2$; (b)~$A$ is in both; (c)~$A$ is
        indexed; (d)~$R_1$ and $R_2$ have the same size.
  \item \textbf{(Short)} Why can a \sqlkw{CHECK} constraint not enforce a
        dependency that spans two tables?
  \item \textbf{(Short)} State the rule for choosing between BCNF and 3NF.
  \item \textbf{(Short)} What three triggers does a summary table need, and
        what goes wrong without the third?
  \item \textbf{(Applied)} Decompose
        \relation{booking}{room, date, slot, course\_code, emp\_id, capacity}
        (\chref{fds}'s review question) into BCNF, test each split for
        losslessness by closure, and say whether every dependency is
        preserved.
  \item \textbf{(Applied)} The university asks for a ``students per
        programme per semester'' figure on its home page, read 5{,}000 times a
        day; about 300 enrolments change a day. Design the summary table and
        its triggers, compute the break-even using this chapter's measured
        costs, and say whether you would build it.
\end{reviewq}
